\documentclass[10pt,twocolumn]{article}

\usepackage[T1]{fontenc}
\usepackage[utf8]{inputenc}
\usepackage{times}
\usepackage[safe]{tipa}
\usepackage[letterpaper,margin=0.75in,columnsep=0.25in]{geometry}
\usepackage{booktabs}
\usepackage{amsmath}
\usepackage{amssymb}
\usepackage[numbers,sort&compress]{natbib}
\usepackage[hidelinks]{hyperref}
\usepackage{url}
\def\UrlBreaks{\do\/\do-}
\usepackage{xcolor}
\usepackage{graphicx}

\setlength{\textfloatsep}{8pt plus 2pt minus 2pt}
\setlength{\floatsep}{6pt plus 2pt minus 2pt}
\setlength{\abovecaptionskip}{4pt}

\newcommand{\att}{Africa TTS}
\newcommand{\attbase}{Africa TTS Base}
\newcommand{\nTrainedLangs}{47}
\newcommand{\nSentences}{2154}
\newcommand{\distN}{45}
\newcommand{\distturbo}{9.1}
\newcommand{\distguided}{9.1}
\newcommand{\distunguided}{10.8}
\newcommand{\turboBeatsMms}{21 of 27}
\newcommand{\turboBeatsOmni}{5 of 29}
\newcommand{\avgMmsturbo}{9.0}
\newcommand{\avgMmsguided}{8.7}
\newcommand{\avgMmsunguided}{10.6}
\newcommand{\avgMmsMms}{10.3}
\newcommand{\avgNMms}{27}
\newcommand{\avgOmniturbo}{8.8}
\newcommand{\avgOmniguided}{8.8}
\newcommand{\avgOmniunguided}{10.4}
\newcommand{\avgOmniOmni}{8.4}
\newcommand{\avgNOmni}{29}
\newcommand{\avgallturbo}{8.6}
\newcommand{\avgallguided}{8.3}
\newcommand{\avgallunguided}{10.5}
\newcommand{\avgallMms}{11.1}
\newcommand{\avgallOmni}{8.2}
\newcommand{\avgNall}{14}
\newcommand{\avgrealreal}{7.5}
\newcommand{\avgrealturbo}{7.7}
\newcommand{\avgrealguided}{7.7}
\newcommand{\avgrealunguided}{8.9}
\newcommand{\avgNreal}{24}
\newcommand{\xrtturbosherpaint}{2.52}
\newcommand{\xrtturbopython}{2.38}
\newcommand{\ttfaturbopython}{0.17}
\newcommand{\xrtfullunguidedsherpaint}{0.82}
\newcommand{\xrtfullunguidedpython}{0.90}
\newcommand{\ttfafullunguidedpython}{0.52}
\newcommand{\xrtfullguidedpython}{0.47}
\newcommand{\ttfafullguidedpython}{0.72}


\title{\att{}: Fast, CPU-Friendly Speech Synthesis\\for 92 African Languages from Phonemes}

\author{%
Mich-Seth Owusu\\[2pt]
\normalsize AfriSpeech\\
\normalsize\texttt{michseth@ghananlp.org}
}

\date{}

\begin{document}
\maketitle

\begin{abstract}
We present \att{}, a 110M-parameter text-to-speech model for 92 African languages and English
that runs faster than real time on an ordinary CPU. All languages share one input space: text is
converted to IPA by rule-based grapheme-to-phoneme tables, so a single model learns sounds rather
than orthographies. \att{} is a 6-layer continuous-latent audio language model built on Pocket
TTS, distilled from a 24-layer base model with classifier-free guidance folded into its weights.
It generates speech \xrtturbopython{}$\times$ faster than real time on four CPU cores, starts
streaming within \ttfaturbopython{}\,s, and fits in a 209\,MB on-device bundle. On held-out
sentences in \nTrainedLangs{} languages, transcribed by Omnilingual ASR, \att{} reaches a mean
character error rate (CER) of \avgMmsturbo{}\% against \avgMmsMms{}\% for MMS-TTS (better in
\turboBeatsMms{} languages), and comes within 0.4 points of the 813M-parameter OmniVoice
(\avgOmniturbo{}\% vs.\ \avgOmniOmni{}\%) at about a seventh of its size. Distillation loses no
intelligibility: \att{} matches its guided base model (\distturbo{}\% vs.\ \distguided{}\%) with
a quarter of the layers and one pass per frame. The model, code, voices and the benchmark are
released.
\end{abstract}

\section{Introduction}

Speech synthesis for African languages is held back less by modelling than by scale and
deployment. Most of the continent's more than two thousand languages have little or no
transcribed speech, and the people who would use a voice interface in those languages often
rely on phones and low-cost laptops without a GPU. Massively multilingual systems have widened
coverage: MMS-TTS~\cite{pratap2023mms} trains one small VITS model~\cite{kim2021vits} per
language for more than a thousand languages, and OmniVoice~\cite{zhu2026omnivoice} clones voices
zero-shot in more than six hundred. But per-language models do not share what languages have in
common, and large zero-shot models need accelerators to run fast.

We present \textbf{\att{}}, a 110M-parameter text-to-speech model for 92 African languages and English,
built on Kyutai's Pocket TTS~\cite{pockettts,rouard2025calm}, a continuous-latent audio language
model designed to run faster than real time on a CPU. Two choices shape the system:
\begin{enumerate}
  \item \textbf{One phonetic input space.} Every language is converted to IPA by rule-based
  grapheme-to-phoneme tables~\cite{africag2p} (espeak-ng~\cite{espeakng} for English), so the
  model learns sounds rather than 92 orthographies, and languages with little data share
  acoustic knowledge with their neighbours.
  \item \textbf{Speech for the CPU.} \att{} has 6 layers (110M parameters including the
  codec). It is distilled from a 24-layer base model, \attbase{}, with classifier-free
  guidance~\cite{ho2022cfg} folded into its weights. It generates speech
  \xrtturbopython{}$\times$ faster than real time on four CPU cores and starts streaming after
  \ttfaturbopython{}\,s. It ships as a 144\,MB 8-bit GGUF model for audio.cpp (desktops and
  servers) and an int8 ONNX bundle for phones.
\end{enumerate}

We evaluate intelligibility with Omnilingual ASR~\cite{omniasr2025} on held-out sentences in
\nTrainedLangs{} languages, against MMS-TTS and OmniVoice, with human recordings of the same sentences as a reference.
Our contributions are: (i) \att{}, a single phoneme-input TTS model covering 92 African
languages, small and fast enough for CPU-only and on-device use, with GGUF and ONNX exports; (ii) a
depth-distillation recipe that keeps the 24-layer base model's guided intelligibility at a
quarter of the layers and one pass per frame; and (iii) an open intelligibility benchmark over \nTrainedLangs{} languages comparing four systems, with all
synthesised audio and transcripts.

The model, voices, code, the benchmark and its outputs are released.\footnote{Code:
\url{https://github.com/AfriSpeech/africa-tts}; model:
\url{https://huggingface.co/AfriSpeech/africa-tts} (base model:
\url{https://huggingface.co/AfriSpeech/africa-tts-base}); live demo:
\url{https://huggingface.co/spaces/AfriSpeech/africa-tts-demo}.}

\section{Related Work}

\paragraph{Multilingual TTS for African languages.} BibleTTS~\cite{meyer2022bibletts} released
studio-quality single-speaker data and VITS baselines for ten sub-Saharan languages. MMS-TTS
\cite{pratap2023mms} scaled per-language VITS models to over 1{,}100 languages from New Testament
recordings, using uroman~\cite{hermjakob2018uroman} romanisation for non-Latin scripts. OmniVoice
\cite{zhu2026omnivoice} is a single 0.6B-class discrete non-autoregressive model trained on
581k hours of open data in over 600 languages, with zero-shot voice cloning. WAXAL
\cite{diack2026waxal} recently added ASR and TTS corpora for over twenty African languages.

\paragraph{Efficient speech language models.} Pocket TTS~\cite{pockettts} follows the
continuous audio language model of \citet{rouard2025calm}: a causal Transformer predicts, frame
by frame at 12.5\,Hz, a continuous latent of the Mimi codec~\cite{defossez2024moshi}, sampled
by a small consistency-trained MLP head, with voice cloning from a short audio prompt. We keep
the architecture, replace the text input with IPA, and distil depth~\cite{hinton2015distilling}.

\paragraph{Evaluation with ASR.} Automatic intelligibility is commonly measured as the word or
character error rate of an ASR system on synthesised speech. Omnilingual ASR~\cite{omniasr2025}
covers over 1{,}600 languages, which for the first time allows a single ASR to score TTS in
most of the languages studied here.

\section{\att{}}

\subsection{Data}
We train on 897 hours of real speech: JW recordings (383\,h), a
crowd-sourced female speech corpus in 17 languages (192\,h), Ghanaian-accented English
(156\,h) and a 200-hour slice of HiFiTTS-2~\cite{langman2025hifitts2} (166\,h). Any clip whose
transcript contains a numeral is dropped (the model reads spelled-out numbers only), as are
clips outside 1--30\,s. Word timings for voice-prompt cut points come from
ctc-forced-aligner~\cite{ctcforcedaligner} for every language (99.95\% of clips aligned). African languages
range from 0.04 to 27 hours each (33 have under two hours); English has 323 hours.

\subsection{Phonetic front end}
Text is converted to IPA word by word: African languages by the C tables of africa-g2p
\cite{africag2p}, English by espeak-ng 1.51 with normalisation. The same C library runs on
device and is byte-identical to the training conversion on every validation sentence. Common
input errors are normalised, e.g.\ Greek and Cyrillic lookalikes of the African letters
\textepsilon{} and \textopeno{}. A SentencePiece unigram tokenizer~\cite{kudo2018sentencepiece}
with 4{,}000 pieces is trained on the IPA, sampling languages in proportion to the square root
of their hours (2.41 tokens per word).

\subsection{Base model}
The 24-layer base model, \attbase{} (336M parameters, of which 20M are the Mimi codec), is
warm-started from Kyutai's English Pocket TTS with a fresh text embedding table for the IPA
vocabulary, and fine-tuned on one H200 GPU with batch 64, learning rate $2\times10^{-4}$ and
20\% text and voice dropout (enabling classifier-free guidance). Training early-stopped at
113k steps; we use step 100k, chosen by listening. It sounds best with guidance (cfg 2.0),
which runs the backbone twice per frame: too slow for a CPU (Table~\ref{tab:speed}).

\subsection{\att{}: depth distillation}
\att{} is a 6-layer student initialised from base-model layers 0, 1, 2, 21, 22 and 23, with the
text table, flow head and end-of-speech head copied, and trained to match the base model's guided
output, so guidance costs nothing at inference. It trains for 200k steps (cosine schedule,
peak $4\times10^{-4}$); among checkpoints at 150k, 175k and 200k, 200k was preferred in
listening tests. Per frame, \att{} computes 6 layers where the guided base model computes 48.

\subsection{Deployment}
\att{} exports to ONNX and runs in sherpa-onnx~\cite{sherpaonnx} with int8 weights in a 209\,MB
bundle that includes the front end and default voices. The exported graphs match PyTorch within $6.5\times10^{-5}$ before quantisation and the decoder at 116\,dB
SNR. For desktops and servers, the model is also converted to an 8-bit GGUF (q8\_0, 144\,MB) and
runs in audio.cpp~\cite{audiocpp}, a C++ engine on ggml, without PyTorch: this is the default engine
of the Python package and of a CPU-only API server (HTTP, WebSocket streaming) that powers a public
demo and scales to zero. The PyTorch model remains available for unquantised and GPU inference.

\section{Evaluation}
\begin{table*}[t]
\centering
\small
\begin{tabular}{lrrrrr}
\toprule
 & \textbf{Real} & \textbf{Africa TTS} (ours) & \textbf{MMS-TTS} & \textbf{OmniVoice} & \textbf{Orpheus} \\
 & \textit{speech} & \textit{110M} & \textit{36M/lang.} & \textit{813M} & \textit{3.3B} \\
\midrule
\multicolumn{6}{l}{\textit{CER (\%)}} \\
Languages MMS-TTS supports (27) &  & \textbf{9.0} & 10.3 &  &  \\
Languages OmniVoice supports (29) &  & 8.8 &  & \textbf{8.4} &  \\
Languages Sunbird Orpheus supports (11) &  & \textbf{5.9} &  &  & 10.3 \\
All systems (6) &  & 6.0 & 7.8 & \textbf{4.3} & 9.8 \\
Languages with real recordings (24) & 7.5 & 7.7 &  &  &  \\
\midrule
\multicolumn{6}{l}{\textit{CER without diacritics (\%)}} \\
Languages MMS-TTS supports (27) &  & \textbf{7.6} & 9.4 &  &  \\
Languages OmniVoice supports (29) &  & 7.5 &  & \textbf{7.0} &  \\
Languages Sunbird Orpheus supports (11) &  & \textbf{4.6} &  &  & 8.8 \\
All systems (6) &  & 4.4 & 6.5 & \textbf{3.0} & 7.7 \\
Languages with real recordings (24) & 6.6 & 6.5 &  &  &  \\
\bottomrule
\end{tabular}
\caption{Intelligibility on languages Africa TTS was trained on: mean per-language character error rate of Omnilingual ASR transcripts (lower is better), averaged over the languages each comparison shares (count in brackets; languages with fewer than 10 test sentences are excluded). Best synthetic system in bold; real speech is the ASR's error on human recordings of the same sentences. Sizes are parameter counts including each system's audio codec or vocoder.}
\label{tab:averages}
\end{table*}


\subsection{Speed}
Table~\ref{tab:speed} reports speed on a dedicated 4-vCPU cloud machine over eight sentences in
eight languages. \att{} runs \xrtturbopython{}$\times$ faster than real time and starts
streaming within \ttfaturbopython{}\,s; the guided base model it was distilled from runs at
\xrtfullguidedpython{}$\times$, slower than real time.

\begin{table}[t]
\centering
\small
\begin{tabular}{llrr}
\toprule
\textbf{Model} & \textbf{Runtime} & \textbf{$\times$RT} $\uparrow$ & \textbf{TTFA} $\downarrow$ \\
\midrule
Africa TTS Turbo & sherpa-onnx int8 & \textbf{2.52} & -- \\
Africa TTS Turbo & PyTorch & 2.38 & 0.17 s \\
Africa TTS unguided & sherpa-onnx int8 & 0.82 & -- \\
Africa TTS unguided & PyTorch & 0.90 & 0.52 s \\
Africa TTS guided & PyTorch & 0.47 & 0.72 s \\
\bottomrule
\end{tabular}
\caption{CPU speed on a dedicated 4-vCPU cloud machine: seconds of audio generated per second of compute ($\times$RT, above 1 is faster than real time) and median time to first audio when streaming (TTFA). PyTorch runs on one thread, sherpa-onnx on four.}
\label{tab:speed}
\end{table}


\subsection{Intelligibility: setup}
\paragraph{Test sentences.} For each trained language we take up to 50 sentences from, in order
of preference, the Omnilingual ASR corpus~\cite{omniasr2025} (test split, else dev, else
train; five languages with tiny splits are topped up from the next split), FLEURS~\cite{conneau2023fleurs} test, WAXAL~\cite{diack2026waxal} test, and FLORES+
\cite{floresplus,nllb2022} devtest (text only). The first three come with human recordings.
Sentences with digits, under 20 or over 220 characters are removed, and Greek or Cyrillic
lookalike letters in FLORES+ are replaced by the intended African letters. None of these
sentences is in our training data.

\paragraph{Systems.} \att{} uses each language's default voice and temperature 0.3, as does
the base model in the distillation comparison. MMS-TTS~\cite{pratap2023mms} is run for every language
with a released model, with uroman for non-Latin scripts. OmniVoice~\cite{zhu2026omnivoice}
is run for every language in its list, cloning the same default voice clip as \att{}; as
OmniVoice also conditions on the prompt's transcript, each prompt is transcribed with the
evaluation ASR where no full human transcript exists. Parameter counts are measured on the
released checkpoints and include each system's codec or vocoder.

\paragraph{Metric.} Each output is transcribed by Omnilingual ASR (\texttt{omniASR\_LLM\_7B\_v2},
conditioned on the language where it has a code), and we report the character error rate (CER)
after Unicode NFC normalisation, lower-casing and removal of punctuation. A second CER also
strips diacritics, as several G2P tables and ASR outputs treat tone marks inconsistently. The
same ASR on the human recordings gives a reference: synthetic speech can score below it, since
read synthetic speech is cleaner than spontaneous or noisy recordings. Means are taken over the
languages each comparison shares.

\subsection{Results}
Table~\ref{tab:averages} summarises intelligibility. Against MMS-TTS, the strongest widely used
per-language baseline, \att{} is more intelligible: it averages \avgMmsturbo{}\% CER against
\avgMmsMms{}\% and is better in \turboBeatsMms{} languages. OmniVoice is the more intelligible
system: on its \avgNOmni{} shared languages it averages \avgOmniOmni{}\% against
\avgOmniturbo{}\% for \att{}, and \att{} is better in only \turboBeatsOmni{} of them, so OmniVoice
wins most languages, by small margins. \att{} closes most of that gap at about a seventh of
OmniVoice's parameters while running faster than real time on a CPU (Table~\ref{tab:speed}).
On the \avgNreal{} languages with human recordings, \att{} (\avgrealturbo{}\%) is within
0.2 points of the ASR's error on real speech (\avgrealreal{}\%). Removing diacritics lowers all
CERs but leaves the ranking unchanged.

\begin{table}[t]
\centering
\footnotesize
\setlength{\tabcolsep}{3.5pt}
\begin{tabular}{lrrrr}
\toprule
 & \textbf{Real} & \textbf{Ours} & \textbf{Base, cfg 2} & \textbf{Base, cfg 1} \\
 & \textit{speech} & \textit{110M} & \textit{336M} & \textit{336M} \\
\midrule
All (45) & & \textbf{9.1} & \textbf{9.1} & 10.8 \\
With recordings (24) & 7.5 & \textbf{7.7} & \textbf{7.7} & 8.9 \\
\bottomrule
\end{tabular}
\caption{Distillation: CER (\%) of \att{} and of the 24-layer base model it was distilled from, with guidance (cfg 2, two passes per frame) or without (cfg 1). Languages with fewer than 10 sentences excluded.}
\label{tab:distill}
\end{table}


\paragraph{Distillation.} Table~\ref{tab:distill} compares \att{} with the base model it was
distilled from. Over all \distN{} languages, \att{} (\distturbo{}\%) matches the guided base
model (\distguided{}\%) and is clearly better than the base model without guidance
(\distunguided{}\%): folding guidance into a 6-layer student keeps its benefit while removing
three quarters of the layers and the second pass, a 5$\times$ speed-up on a CPU.

\section{Limitations}
The ASR-based metric measures intelligibility only, not naturalness or speaker similarity, and
inherits the ASR's own errors; for languages whose test text comes from the Omnilingual ASR
corpus's training split, the ASR may have seen those sentences, which favours all systems
equally. FLEURS sentences are FLORES sentences read aloud. Three languages (Chitumbuka, Swati,
Kimbundu) have no ASR language code and are transcribed unconditioned. Most of our African
training speech is read religious text, so conversational and technical registers are less
natural, and languages with under about two hours of data are weak. Some G2P tables do not
mark tone. Voice cloning can be
misused to impersonate people; the models are released for non-commercial use under
CC-BY-NC-4.0.

\section{Conclusion}
\att{} shows that one 110M-parameter phoneme-input model can serve 92 African languages
faster than real time on a CPU: distilled from a 24-layer base model, it keeps that model's
guided intelligibility. It
is clearly more intelligible than per-language MMS-TTS models, and close to OmniVoice at a
fraction of its size. Future work includes
human listening tests with native speakers, more conversational training speech, and tone
marking in more G2P tables.

\bibliographystyle{plainnat}
\bibliography{refs}

\onecolumn
\appendix
\section{Per-language results}
\begin{table}[h!]
\centering
\small
\setlength{\tabcolsep}{4pt}
\begin{tabular}{llcrrrrr}
\toprule
\textbf{Language} & \textbf{Code} & \textbf{Src} & \textbf{$n$} & \textbf{Real} & \textbf{Africa TTS} (ours) & \textbf{MMS-TTS} & \textbf{OmniVoice} \\
 &  &  &  & \textit{speech} & \textit{110M} & \textit{36M/lang.} & \textit{813M} \\
\midrule
Amharic & \texttt{amh} & FL & 50 & 8.0 & 5.7 & 10.9 & \textbf{4.8} \\
Changana (Mozambique) & \texttt{tso} & FR & 50 & -- & 11.5 & \textbf{5.0} & -- \\
Chitumbuka$^\dagger$ & \texttt{tum} & FR & 50 & -- & 11.5 & -- & -- \\
Chokwe & \texttt{cjk} & FR & 50 & -- & 14.9 & -- & \textbf{12.9} \\
Cibemba & \texttt{bem} & FR & 50 & -- & \textbf{5.3} & 7.3 & -- \\
Cinyanja & \texttt{nya} & FL & 50 & 8.0 & 5.0 & 7.1 & \textbf{2.6} \\
Dagaare & \texttt{dga} & WX & 50 & 22.1 & \textbf{10.4} & 10.6 & -- \\
English & \texttt{eng} & FL & 50 & 2.3 & 2.4 & 3.5 & \textbf{2.0} \\
Ewe & \texttt{ewe} & FR & 50 & -- & 6.4 & \textbf{6.2} & -- \\
Fante & \texttt{fat} & OC & 50 & 13.9 & \textbf{7.3} & -- & 8.6 \\
Fon & \texttt{fon} & FR & 50 & -- & 22.5 & \textbf{20.9} & -- \\
Frafra & \texttt{gur} & OC & 50 & 5.5 & 9.8 & -- & \textbf{9.1} \\
Hausa & \texttt{hau} & FL & 50 & 4.7 & 4.2 & 6.2 & \textbf{3.6} \\
Herero & \texttt{her} & OC & 14 & 3.8 & 4.8 & -- & \textbf{1.5} \\
Igbo & \texttt{ibo} & FL & 50 & 8.2 & 6.9 & -- & \textbf{6.6} \\
Isoko & \texttt{iso} & OC & 5 & 13.8 & 29.4 & -- & \textbf{14.9} \\
Jula & \texttt{dyu} & FR & 50 & -- & \textbf{23.8} & 29.7 & 47.9 \\
Kabiye & \texttt{kbp} & FR & 50 & -- & \textbf{11.1} & 11.7 & -- \\
Kabuverdianu & \texttt{kea} & FL & 50 & 2.6 & 11.9 & -- & \textbf{8.6} \\
Kabyle & \texttt{kab} & FR & 50 & -- & 14.0 & 14.5 & \textbf{6.8} \\
Kamba & \texttt{kam} & FL & 50 & 10.3 & 14.6 & -- & \textbf{13.5} \\
Kikuyu & \texttt{kik} & FR & 50 & -- & \textbf{7.0} & 8.3 & -- \\
Kimbundu$^\dagger$ & \texttt{kmb} & FR & 50 & -- & 12.5 & -- & -- \\
Kinyarwanda & \texttt{kin} & FR & 50 & -- & 3.2 & 6.3 & \textbf{1.8} \\
Kirundi & \texttt{run} & FR & 50 & -- & \textbf{3.0} & 10.2 & -- \\
Kpelle & \texttt{xpe} & OC & 1 & 21.1 & 34.1 & -- & \textbf{26.0} \\
Lingala & \texttt{lin} & FL & 50 & 3.1 & 3.8 & -- & \textbf{3.4} \\
Luganda & \texttt{lug} & FL & 50 & 7.3 & 8.2 & 8.8 & \textbf{5.2} \\
Macua & \texttt{vmw} & FR & 50 & -- & 6.7 & \textbf{6.0} & -- \\
Malagasy & \texttt{mlg} & OC & 50 & 3.6 & 8.8 & 9.6 & \textbf{5.2} \\
Mauritian Creole & \texttt{mfe} & FR & 50 & -- & 7.1 & \textbf{5.9} & -- \\
Meru & \texttt{mer} & OC & 21 & 9.2 & \textbf{14.1} & -- & 25.6 \\
Mooré & \texttt{mos} & FR & 50 & -- & 15.4 & \textbf{14.1} & -- \\
Ndonga & \texttt{ndo} & OC & 13 & 14.0 & 5.9 & -- & \textbf{5.6} \\
Sango & \texttt{sag} & FR & 50 & -- & \textbf{10.0} & 10.6 & -- \\
Sena & \texttt{seh} & FR & 50 & -- & \textbf{6.2} & 7.3 & -- \\
Sepedi & \texttt{nso} & FL & 50 & 5.9 & 4.9 & -- & \textbf{4.1} \\
Shona & \texttt{sna} & FL & 50 & 2.9 & 3.0 & 6.1 & \textbf{1.9} \\
Swahili & \texttt{swh} & FL & 50 & 1.8 & 1.2 & 2.6 & \textbf{0.6} \\
Swati$^\dagger$ & \texttt{ssw} & FR & 50 & -- & 9.5 & -- & -- \\
Tigrinya & \texttt{tir} & FR & 50 & -- & 11.9 & 17.6 & \textbf{7.7} \\
Twi & \texttt{twi} & FR & 50 & -- & \textbf{12.1} & 13.1 & 12.5 \\
Umbundu & \texttt{umb} & FL & 50 & 9.5 & 16.6 & -- & \textbf{14.8} \\
Wolof & \texttt{wol} & FL & 50 & 10.3 & 12.2 & -- & \textbf{7.2} \\
Xhosa & \texttt{xho} & FL & 50 & 3.4 & 4.3 & -- & \textbf{3.5} \\
Yoruba & \texttt{yor} & FL & 50 & 15.9 & 17.0 & 19.3 & \textbf{12.8} \\
Zulu & \texttt{zul} & FL & 50 & 3.7 & \textbf{2.2} & -- & 2.4 \\
\bottomrule
\end{tabular}
\caption{CER (\%) per trained language. Src: OC Omnilingual ASR corpus, FL FLEURS, WX WAXAL, FR FLORES+ (text only, no real recordings). $n$: sentences. $^\dagger$: no Omnilingual ASR language code; transcribed without one. Best synthetic system in bold.}
\label{tab:perlang}
\end{table}


\end{document}
